\section{A log-concave cumulative reverse profile}

For a host $Z$ of weight $c$, put
\[
 q_Z(z)=z^cF_Z(1/z),\qquad
 A_Z(k)=\sum_{j=0}^k[z^j]q_Z(z),
\]
with $A_Z(k)=0$ for $k<0$. The cumulative sequence is constant for $k\ge c$.
Write $[t]=1+z+\cdots+z^{t-1}$ for $t\ge1$. For a fixed
legal peak $M\ge\max Z$, set
\[
 D_Z=1+(S_M-1)F_Z,\qquad
 \delta_Z(z)=z^{c+M}D_Z(1/z).
\]
The kernel is nonnegative, since $D_Z=F_{ZM}-F_Z$.
Parenthesized integer arguments of $D_Z$ and $\delta_Z$ below denote
coefficients, not polynomial evaluations. We first prove the two
cumulative properties when $Z$ is increasing.

\begin{theorem}
For every increasing arm $Z$, the sequence $A_Z$ is log-concave:
\[
 A_Z(k)^2\ge A_Z(k-1)A_Z(k+1)\qquad(k\in\mathbb Z).
\]
In particular, for a legal peak $M$, the reversed kernel $\delta_Z$ is
log-concave at each index $0\le k\le M-2$.
\end{theorem}

\begin{proof}
The empty arm gives $A_\varnothing(k)=\mathbf1_{k\ge0}$, which is
log-concave with interval support. Suppose the result holds for an arm of
weight $c$ and append a global maximum $t\ge1$. The reversed recurrence is
\[
 q_{Zt}=[t]q_Z+z^{c+t}.
\]
Consequently
\[
 A_{Zt}(k)=W(k)+\mathbf1_{k\ge c+t},\qquad
 W(k)=\sum_{i=0}^{t-1}A_Z(k-i).
\]
The convolution of two nonnegative log-concave sequences with interval
support is log-concave; it applies here to $A_Z$ and $[t]$. Although $A_Z$
has infinite constant support to the right, this application is local:
truncate it strictly after all indices occurring in any of the three
coefficients being compared. Such a terminal truncation preserves
log-concavity and interval support.

Set $C=tA_Z(c)$. Since $[z^c]q_Z=[z^0]F_Z=1$, the last two values before
the added step are
\[
 W(c+t-2)=C-1,\qquad W(c+t-1)=C,
\]
and $W(k)=C$ thereafter. These formulas also cover $c=0$ and $t=1$, with
zero extension. All log-concavity inequalities before index $c+t-1$
remain unchanged. The boundary values are $C-1,C,C+1,C+1,\ldots$.
The only two remaining boundary inequalities are
\[
 C^2-(C-1)(C+1)=1>0,\qquad
 (C+1)^2-C(C+1)=C+1>0.
\]
Afterward the sequence is constant. Induction proves the first assertion.

The kernel identity is
\[
 \delta_Z=([M]-z^M)q_Z+z^{c+M}.
\]
Thus $[z^k]\delta_Z=A_Z(k)$ for $0\le k<M$. All three coefficients in a
log-concavity comparison agree when $0\le k\le M-2$, including the zero
coefficient at index $-1$. This proves the second assertion. No extension
past that interval is asserted.
\end{proof}

\begin{theorem}[Cumulative reverse kernel]
Fix $M\ge2$. For every increasing arm $Y$ with parts strictly below $M$,
the sequence
\[
 B_Y(k)=\sum_{j=0}^k[z^j]\delta_Y(z),\qquad B_Y(k)=0\ (k<0),
\]
is log-concave on all integer indices.
\end{theorem}

\begin{proof}
The empty arm gives $\delta_\varnothing=[M]$, so the assertion is clear.
Suppose $Y=Zt$, put $c=\wt Z$, and let $b=c+t$. The reversed recurrence
gives
\[
 B_Y(k)=W(k)+\sum_{j=0}^k[z^j]\bigl(z^b[M-t]\bigr),\qquad
 W=[t]*B_Z.
\]
Induction and the log-concave convolution lemma make $W$ log-concave.
For a comparison centered at $k\le b-2$, the correction has not started,
and all three coefficients agree with $W$.

For $k\ge b$, the kernel's known low-rank growth says
$\delta_Y(k+1)\le\delta_Y(k)$: after reversing, these are forward ranks
at most $M$. Consequently $B_Y$ is nonnegative, nondecreasing and
concave at such a comparison, which implies log-concavity there.
Indeed, if consecutive increments are $s\ge t\ge0$ and the first value
is $A\ge0$, the log-concavity determinant is $s^2+A(s-t)\ge0$.

Only $k=b-1$ remains. It is covered by the same argument if
$\delta_Y(b)\le\delta_Y(b-1)$. Otherwise the reversed recurrence gives
\[
 0<\delta_Y(b)-\delta_Y(b-1)
   =1+D_Z(M-t)-D_Z(M).
\]
The two forward coefficients are integers and $D_Z$ is nondecreasing
through $M$, so they must be equal. In the kernel proof, the identity
\[
 (1-z)D_Z=E_Z-z^{M+1}F_Z
\]
and the positive interval support $\supp E_Z=[q,c+1]$, with $q\le t+1$,
then force $c+1\le M-t$: otherwise the interval $[M-t+1,M]$ meets the
positive support and $D_Z(M)>D_Z(M-t)$. Hence $b<M$.

For every $0\le j\le b$, therefore, $\delta_Y(j)=A_Y(j)$ by the preceding
theorem. The three values of $B_Y$ at the remaining comparison are
cumulative sums of the log-concave sequence $A_Y$. Cumulative summation
preserves log-concavity, by convolution with the sequence
$\mathbf1_{k\ge0}$ (or a sufficiently long finite window). This proves
the final comparison and completes the induction.
\end{proof}

